% !Mode:: "TeX:UTF-8"
% !TEX program = xelatex

% =============================================
% 字体定义与大小设置
% =============================================
\newcommand{\song}{\songti}     % 宋体
\newcommand{\fs}{\fangsong}     % 仿宋体
\newcommand{\kai}{\kaishu}      % 楷体
\newcommand{\hei}{\heiti}       % 黑体
\newcommand{\li}{\lishu}        % 隶书

% --- 全局英文字体（Windows 优先，Linux 回退 Nimbus / gsfonts）---
\IfFontExistsTF{Times New Roman}
  {\setmainfont{Times New Roman}}
  {\setmainfont{Nimbus Roman}}
\IfFontExistsTF{Arial}
  {\setsansfont{Arial}}
  {\setsansfont{Nimbus Sans}}
\IfFontExistsTF{Courier New}
  {\setmonofont{Courier New}}
  {\setmonofont{Nimbus Mono PS}}

\newcommand{\xiaochu}{\fontsize{36pt}{30pt}\selectfont}     % 小初
\newcommand{\yihao}{\fontsize{28pt}{26pt}\selectfont}       % 一号
\newcommand{\xiaoyi}{\fontsize{24pt}{24pt}\selectfont}      % 小一
\newcommand{\erhao}{\fontsize{21pt}{1.25\baselineskip}\selectfont} % 二号
\newcommand{\xiaoer}{\fontsize{18pt}{18pt}\selectfont}      % 小二
\newcommand{\sanhao}{\fontsize{15.75pt}{15.75pt}\selectfont} % 三号
\newcommand{\xiaosan}{\fontsize{15pt}{15pt}\selectfont}     % 小三
\newcommand{\sihao}{\fontsize{14pt}{14pt}\selectfont}       % 四号
\newcommand{\xiaosi}{\fontsize{12pt}{12pt}\selectfont}      % 小四
\newcommand{\wuhao}{\fontsize{10.5pt}{10.5pt}\selectfont}   % 五号
\newcommand{\xiaowu}{\fontsize{9pt}{9pt}\selectfont}        % 小五

% =============================================
% 全局段落与列表设置
% =============================================
\setitemize{leftmargin=3em,itemsep=0em,partopsep=0em,parsep=0em,topsep=-0em}
\setenumerate{leftmargin=3em,itemsep=0em,partopsep=0em,parsep=0em,topsep=0em}
\setlength{\parindent}{2em}

\makeatletter
\renewcommand\normalsize{
  \@setfontsize\normalsize{12pt}{12pt} % 小四对应 12 pt
  \setlength\abovedisplayskip{4pt}
  \setlength\abovedisplayshortskip{4pt}
  \setlength\belowdisplayskip{\abovedisplayskip}
  \setlength\belowdisplayshortskip{\abovedisplayshortskip}
  \let\@listi\@listI
  \setlength{\baselineskip}{20pt}
}
\def\defaultfont{\renewcommand{\baselinestretch}{1.5}\normalsize\selectfont}
\renewcommand{\CJKglue}{\hspace -0.1 pt plus 0.08\baselineskip}
\makeatother

% =============================================
% 章节标题格式 (统一使用 ctexset)
% =============================================
\setcounter{secnumdepth}{4}

\ctexset{
    chapter = {
        name = {第,章},
        format = \linespread{1.0}\zihao{-3}\heiti\centering, % 恢复为主文件要求的字号
        number = \arabic{chapter},
        aftername = \quad,
        beforeskip = {7pt},
        afterskip = {18pt},
        pagestyle = plain,
    },
    section={
        format = \zihao{-3}\heiti\raggedright,
        name = {},
        number = \arabic{chapter}.\arabic{section},
        beforeskip = 1.0ex plus 0.2ex minus .2ex,
        afterskip = 1.0ex plus 0.2ex minus .2ex,
        aftername = \quad
    },
    subsection={
        format = \zihao{-4}\heiti\raggedright,
        name = {},
        number = \arabic{chapter}.\arabic{section}.\arabic{subsection},
        beforeskip = 1.0ex plus 0.2ex minus .2ex,
        afterskip = 1.0ex plus 0.2ex minus .2ex,
        aftername = \quad
    },
    %subsubsection={
    %    format = \zihao{-4}\fangsong\centering,
    %    number = \ttfamily\arabic{subsubsection},
    %    name = {,.},
    %    beforeskip = 1.0ex plus 0.2ex minus .2ex,
    %    afterskip = 1.0ex plus 0.2ex minus .2ex,
    %    aftername = \hspace{0pt}
    %}
}

% =============================================
% 图表与公式编号规范
% =============================================
\renewcommand{\tablename}{表}                                     
\renewcommand{\figurename}{图}                                    
\renewcommand{\thefigure}{\arabic{chapter}-\arabic{figure}}       
\renewcommand{\thetable}{\arabic{chapter}-\arabic{table}}         
\renewcommand{\theequation}{\arabic{chapter}-\arabic{equation}}   
\makeatletter
\renewcommand{\tagform@}[1]{\maketag@@@{式(#1\@@italiccorr)}}
\makeatother
\renewcommand{\eqref}[1]{\textup{式(\ref{#1})}}
\newcommand{\ud}{\mathrm{d}}

% 浮动体标题样式
\makeatletter
\long\def\@makecaption#1#2{
   \vskip\abovecaptionskip
   \sbox\@tempboxa{\centering\wuhao\song{#1\quad #2} }
   \ifdim \wd\@tempboxa >\hsize
     \centering\wuhao\song{#1\quad #2} \par
   \else
     \global \@minipagefalse
     \hb@xt@\hsize{\hfil\box\@tempboxa\hfil}
   \fi
   \vskip\belowcaptionskip}
\makeatother

% =============================================
% 自定义命令与环境
% =============================================
% 表格单元格内换行
\newcommand{\tabincell}[2]{\begin{tabular}{@{}#1@{}}#2\end{tabular}}
% 表格列格式扩展
\newcommand{\PreserveBackslash}[1]{\let\temp=\\#1\let\\=\temp}
\newcolumntype{C}[1]{>{\PreserveBackslash\centering}p{#1}}
\newcolumntype{R}[1]{>{\PreserveBackslash\raggedleft}p{#1}}
\newcolumntype{L}[1]{>{\PreserveBackslash\raggedright}p{#1}}

% 避免 hyperref 和部分包冲突导致目录链接失效
\def\temp{\relax}
\let\temp\addcontentsline
\gdef\addcontentsline{\phantomsection\temp}

% 代码高亮由 setup/minted-setup.tex（minted + Pygments）统一配置

% 定理环境
\theoremstyle{plain}
\theorembodyfont{\song\rmfamily}
\theoremheaderfont{\hei\rmfamily}
\newtheorem{theorem}{定理~}[chapter]
\newtheorem{lemma}{引理~}[chapter]
\newtheorem{definition}{定义~}[chapter]
\newtheorem{example}{例~}[chapter]
\newenvironment{proof}{\noindent{\hei 证明：}}{\hfill $ \square $ \vskip 4mm}
\theoremsymbol{$\square$}

% =============================================
% 参考文献排版设置
% =============================================
\renewcommand{\bibname}{参考文献}

% 1. 强制 natbib 全局使用带方括号的上标格式（如：^[1]）
\setcitestyle{super,square}

% 2. 贴心预留：如果你在正文中需要说“请参见文献[1]”（不作为上标）
% 以后遇到这种情况，请使用 \inlinecite{...} 命令
\newcommand{\inlinecite}[1]{[\citenum{#1}]}

\addtolength{\bibsep}{-0.5em} 
\setlength{\bibhang}{2em}     

% =============================================
% 目录样式控制 (titletoc)
% =============================================
\renewcommand{\contentsname}{\xiaoer 目\qquad 录}
\setcounter{tocdepth}{1} 

\titlecontents{chapter}[0em]{\vspace{2pt}\sihao\song\bfseries}
             {\thecontentslabel\quad}{\mbox{}} 
             {\hspace{.5em}\titlerule*[5pt]{$\cdot$}\sihao\contentspage}

%\titlecontents{chapter}[0em]{\vspace{2pt}\sihao\song\bfseries}
%             {\thecontentslabel \quad}{}
%             {\hspace{.5em}\titlerule*[5pt]{$\cdot$}\sihao\contentspage}

\titlecontents{section}[2em]{\vspace{2pt}\xiaosi\song}
             {\thecontentslabel\quad}{}
             {\hspace{.5em}\titlerule*[3pt]{$\cdot$}\xiaosi\contentspage}
             
\titlecontents{subsection}[4em]{\vspace{2pt}\xiaosi\song}
             {\thecontentslabel 、}{}
             {\hspace{.5em}\titlerule*[3pt]{$\cdot$}\xiaosi\contentspage}

% =============================================
% 页眉页脚与封面定义
% =============================================
\usepackage{fancyhdr}
\makeatletter
\pagestyle{fancy}
\fancyhf{}
\fancyhead[C]{\song\wuhao \@cheading} 
\fancyfoot[C]{\song\xiaowu ~\thepage~}

\def\ctitle#1{\def\@ctitle{#1}}\def\@ctitle{}
\def\cntitle#1{\def\@cntitle{#1}}\def\@cntitle{}
\def\cctitle#1{\def\@cctitle{#1}}\def\@cctitle{}
\def\etitle#1{\def\@etitle{#1}}\def\@etitle{}
\def\cauthor#1{\def\@cauthor{#1}}\def\@cauthor{}
\def\csupervisor#1{\def\@csupervisor{#1}}\def\@csupervisor{}
\def\caffil#1{\def\@caffil{#1}}\def\@caffil{}
\def\csubject#1{\def\@csubject{#1}}\def\@csubject{}
\def\csubjectcode#1{\def\@csubjectcode{#1}}\def\@csubjectcode{}
\def\cresearch#1{\def\@cresearch{#1}}\def\@cresearch{}
\def\cgrade#1{\def\@cgrade{#1}}\def\@cgrade{}
\def\cnumber#1{\def\@cnumber{#1}}\def\@cnumber{}
\def\cheading#1{\def\@cheading{#1}}\def\@cheading{}
\long\def\cabstract#1{\long\def\@cabstract{#1}}\long\def\@cabstract{}
\long\def\eabstract#1{\long\def\@eabstract{#1}}\long\def\@eabstract{}
\def\ckeywords#1{\def\@ckeywords{#1}}\def\@ckeywords{}
\def\ekeywords#1{\def\@ekeywords{#1}}\def\@ekeywords{}


\makeatother

% 强制取消 \emph 命令产生的下划线，将其恢复为标准的斜体
% 彻底重置 \emph 命令，使其在任何情况下都保持正文默认字体（即：不倾斜、不加下划线）
\let\oldemph\emph
\renewcommand{\emph}[1]{#1}